\documentclass[a4paper,12pt]{article} % {{{ Head

% -------------------------------------------------
%
% | |    ___   / ___| |__   ___ _ __
% | |   / _ \ | |   | '_ \ / _ \ '_ \
% | |__|  __/ | |___| | | |  __/ | | |
% |_____\___|  \____|_| |_|\___|_| |_|
%
%  Created on Sat 17 Jul 2021 03:50:06 PM EDT
%
% -------------------------------------------------
% Set up the page margins.
% -------------------------------------------------
% \usepackage[text={6.5in,9.5in},centering]{geometry}
\usepackage[margin=1in]{geometry}
% \setlength{\topmargin}{-0.25in}

% -------------------------------------------------
% Load Packages here
% -------------------------------------------------
\usepackage{graphicx,url,subfig}
% \RequirePackage[OT1]{fontenc}
% \usepackage{etex}
\RequirePackage{amsthm,amsmath,amssymb,amscd,mathrsfs}
\RequirePackage[colorlinks,citecolor=blue,urlcolor=blue]{hyperref}
\RequirePackage[numbers]{natbib}
% \usepackage{luacode}
\usepackage{enumerate}
\usepackage{datetime}
\usepackage[normalem]{ulem} % either use this (simple) or
\usepackage{soul} % use this (many fancier options)
\usepackage{xcolor}

% -------------------------------------------------
% check references
% -------------------------------------------------
% \usepackage{refcheck}
% \usepackage[notref,notcite]{showkeys}
% \usepackage[notcite]{showkeys}
% \usepackage{showkeys}

% -------------------------------------------------
% Environments here
% -------------------------------------------------
\theoremstyle{plain}
\newtheorem{theorem}{Theorem}[section]
\newtheorem{corollary}[theorem]{Corollary}
\newtheorem{lemma}[theorem]{Lemma}
\newtheorem{proposition}[theorem]{Proposition}

\theoremstyle{definition}
\newtheorem{definition}[theorem]{Definition}
\newtheorem{hypothesis}[theorem]{Hypothesis}
\newtheorem{assumption}[theorem]{Assumption}

\newtheorem{remark}[theorem]{Remark}
\newtheorem{conjecture}[theorem]{Conjecture}
\newtheorem{example}[theorem]{Example}

% -------------------------------------------------
%  Define short hand symbols.
% -------------------------------------------------
\newcommand{\E}{\mathbb{E}}
\newcommand{\W}{\dot{W}}
\newcommand{\ud}{\ensuremath{\mathrm{d} }}
\newcommand{\Ceil}[1]{\left\lceil #1 \right\rceil}
\newcommand{\Floor}[1]{\left\lfloor #1 \right\rfloor}
\newcommand{\sgn}{\text{sgn}}
\newcommand{\Norm}[1]{\left|\left|  #1   \right|\right|}
\newcommand{\Ito}{It\^{o} }
\newcommand{\Itos}{It\^{o}'s }
\newcommand{\spt}[1]{\text{supp}\left(#1\right)}
\newcommand{\InPrd}[1]{\left\langle #1 \right\rangle}
\newcommand{\arctanh}{\operatorname{arctanh}}
\DeclareMathOperator{\esssup}{\ensuremath{ess\,sup}}
\newcommand{\calB}{\mathcal{B}}
\newcommand{\calD}{\mathcal{D}}
\newcommand{\calE}{\mathcal{E}}
\newcommand{\calF}{\mathcal{F}}
\newcommand{\calG}{\mathcal{G}}
\newcommand{\calK}{\mathcal{K}}
\newcommand{\calH}{\mathcal{H}}
\newcommand{\calI}{\mathcal{I}}
\newcommand{\calL}{\mathcal{L}}
\newcommand{\calM}{\mathcal{M}}
\newcommand{\calN}{\mathcal{N}}
\newcommand{\calO}{\mathcal{O}}
\newcommand{\calT}{\mathcal{T}}
\newcommand{\calP}{\mathcal{P}}
\newcommand{\calR}{\mathcal{R}}
\newcommand{\calS}{\mathcal{S}}
\newcommand{\bbC}{\mathbb{C}}
\newcommand{\bbN}{\mathbb{N}}
\newcommand{\bbZ}{\mathbb{Z}}
\newcommand*{\one}{{ {\rm 1\mkern-1.5mu}\!{\rm I} }}
\newcommand{\R}{\mathbb{R}}
\newcommand{\RR}{\mathbb{R}}
\newcommand{\myEnd}{\hfill$\square$}
\newcommand{\Erf}{\ensuremath{\mathrm{erf} }}
\newcommand{\Erfi}{\ensuremath{\mathrm{erfi} }}
\newcommand{\Erfc}{\ensuremath{\mathrm{erfc} }}
\newcommand{\He}{\ensuremath{\mathrm{He} }}
\DeclareMathOperator{\Lip}{Lip}
\DeclareMathOperator{\LIP}{Lip}
\DeclareMathOperator{\lip}{\mathit{l}}
\newcommand{\Dxa}{{}_x D_\delta^a}
\newcommand{\Dtb}{{}_t D_*^\beta}
\newcommand{\lMr}[3]{\:{}_{#1} #2_{#3}}
\newcommand{\tlMr}[4]{\:{}^{\hspace{0.2em}#1}_{#2} \hspace{-0.1em}#3_{#4}}
% }}}

% -------------------------------------------------
% Title, authors and date go here.
% -------------------------------------------------
\begin{document}

% -------------------------------------------------
% Document begins here
% -------------------------------------------------

Here let's mimic the proof of \cite[Lemma 7.2]{ChenLi04}.
\bigskip

By Lagrange multiplier, we see that
\begin{multline*}
	2\left(\int_\R |f(x)|^{2p}\ud x\right)^{-(p-1)/p} \int_\R f^{2p-1}(x)g(x)\ud x
	- \frac{1}{2\pi} \int_{\R} |\xi|^a \mathcal{F} g(\xi) \overline{\mathcal{F} f(\xi)} \ud \xi \\
	= \frac{\lambda}{\pi}\int_\R \overline{\mathcal{F}f(\xi)} \mathcal{F}g(\xi) \ud \xi.
\end{multline*}
We can remove the conjugate sign in the above equality because $f$ is an even function.
By Plancherel theorem, we see that
\begin{multline*}
	\left(\int_\R |f(x)|^{2p}\ud x\right)^{-(p-1)/p} \int_\R \mathcal{F}[f^{2p-1}](\xi)
	\mathcal{F}g(\xi)\ud \xi -
	\frac{1}{2}\int_{\R} |\xi|^a \mathcal{F} g(\xi) \mathcal{F} f(\xi) \ud \xi \\
	= \lambda\int_\R \mathcal{F}f(\xi) \mathcal{F}g(\xi) \ud \xi.
\end{multline*}
By writing the above equality in the following way:
\begin{multline*}
	\left(\int_\R |f(x)|^{2p}\ud x\right)^{-(p-1)/p} \int_\R
	\frac{\mathcal{F}[f^{2p-1}](\xi)}{|\xi|^{a/2}}
	\mathcal{F}g(\xi)|\xi|^{a/2}\ud \xi -
	\frac{1}{2}\int_{\R} \mathcal{F} f(\xi) |\xi|^{a/2} \mathcal{F} g(\xi) |\xi|^{a/2} \ud \xi \\
	= \lambda\int_\R \frac{\mathcal{F}f(\xi)}{|\xi|^{a/2}} \mathcal{F}g(\xi) |\xi|^{a/2} \ud \xi,
\end{multline*}
we see that for all $\xi\in\R$,
\begin{align*}
	\left(\int_\R |f(x)|^{2p}\ud x\right)^{-(p-1)/p} \frac{\mathcal{F}[f^{2p-1}](\xi)}{|\xi|^{a/2}}
	- \frac{1}{2}\mathcal{F} f(\xi) |\xi|^{a/2}
	= \lambda \frac{\mathcal{F}f(\xi)}{|\xi|^{a/2}}.
\end{align*}
Multiply $|\xi|^{a/2}$ on both sides to see that
\begin{align}
	\label{E:(7.9)} \tag{7.9}
	\left(\int_\R |f(x)|^{2p}\ud x\right)^{-(p-1)/p} \mathcal{F}[f^{2p-1}](\xi)
	- \frac{1}{2}\mathcal{F} f(\xi) |\xi|^{a}
	= \lambda \mathcal{F}f(\xi).
\end{align}
Now multiply $(2\pi)^{-1}\mathcal{F}f(\xi)$ on both sides of \eqref{E:(7.9)} and integrate to see that
\begin{align*}
	\left(\int_\R |f(x)|^{2p}\ud x\right)^{1/p}
	-\frac{1}{4\pi}\int_\R |\xi|^a |\mathcal{F}f(\xi)|^2\ud \xi
	= \lambda,
\end{align*}
where we have used the fact that $\Norm{f}_{L^2(\R)}=1$.

Now multiply $(2\pi)^{-1}\mathcal{F}f(\xi)|\xi|^{a/2}$ on both sides of \eqref{E:(7.9)} and
integrate to see that... \textcolor{magenta}{The difficulties come from here. I cannot find the
counter part of the integration by parts for fractional derivatives.}


% {{{ References
% -------------------------------------------------
% References go here.
% -------------------------------------------------
\begin{thebibliography}{999}

	\bibitem{ChenLi04} % {{{ 11.
	Chen, X. and Li, W. (2004).
	\newblock Large and moderate deviations for intersection local times.
	\newblock {\em Probab. Theory Related Fields} {\bf 128} 213--254.
	% }}}

\end{thebibliography} % }}}

\end{document}

